Neural computers benefit from progress across the basic pillars while gradually enabling CNC extensions; these capabilities co-evolve rather than forming a strict prerequisite chain. We identify two complementary evolution paths.

\noindent\textbf{Consolidating the basic pillars.} Current pixel-grounded approaches establish early I/O \& control, but significant gaps remain. Efficiency demands streaming inference under realistic compute budgets and latency constraints. Memory \& storage needs auditable, long-term information that can be accessed through the latent runtime state. Computation \& reasoning requires dynamics that support multi-step reasoning and verifiable checks. Hybrid architectures combining video-based latent-state models with structured renderers can address memory and efficiency bottlenecks while preserving pixel-grounded perception.

\noindent\textbf{Enabling CNC extensions through neural artifacts.} As the pillars become more reliable, CNC extensions become more achievable through neural networks generating downstream neural modules. Tool bridges can emerge when NCs synthesize interface adapters. Condition-driven generalization becomes possible when conditions act as specifications and latent-state dynamics can reliably switch modes without human-written bespoke programs. Programmability requires persistent neural routines (``latent programs'') that can be stored and invoked. Artifact generation represents the ultimate capability: NCs that emit executable modules extending their own capabilities.

This suggests a concrete research agenda: near-term emphasis on execution efficiency and memory/storage through hybrid architectures; mid-term efforts on computation \& reasoning and robust I/O/control; and long-term pursuit of programmability and artifact generation for truly self-extending neural computers.
